\subsection{Parallel Lines and Proportional Segments}

A central fact of Euclidean geometry is that parallel lines preserve ratios.
This is the content of Thales' theorem in the form needed here.

\begin{theorembox}
Let $D$ lie on $AB$ and $E$ lie on $AC$ in a triangle $ABC$. If
\[
DE\parallel BC,
\]
then
\[
\frac{AD}{AB}
=
\frac{AE}{AC}
=
\frac{DE}{BC}.
\]
\end{theorembox}

\begin{center}
\begin{tikzpicture}[scale=0.9]
    \coordinate (A) at (0,3.6);
    \coordinate (B) at (-3.2,0);
    \coordinate (C) at (3.8,0);
    \coordinate (D) at (-1.6,1.8);
    \coordinate (E) at (1.9,1.8);
    \draw[subset] (A) -- (B) -- (C) -- cycle;
    \draw[very thick] (D) -- (E);
    \fill[geometricSubsetColor] (A) circle (2pt) node[above] {$A$};
    \fill[geometricSubsetColor] (B) circle (2pt) node[below left] {$B$};
    \fill[geometricSubsetColor] (C) circle (2pt) node[below right] {$C$};
    \fill[geometricSubsetColor] (D) circle (2pt) node[left] {$D$};
    \fill[geometricSubsetColor] (E) circle (2pt) node[right] {$E$};
    \node at (0.15,2.05) {$DE\parallel BC$};
\end{tikzpicture}
\end{center}

The smaller triangle $ADE$ and the larger triangle $ABC$ have the same three
angles. Their corresponding sides differ by one common scale factor. This is
the basic phenomenon of similarity.

\begin{definitionbox}
Two triangles are similar when their corresponding angles are equal. In that
case their corresponding side lengths are proportional.
\end{definitionbox}

Similarity lets us change the size of a figure while preserving its shape. It
is the reason ratios of corresponding lengths can depend only on angles.
